\documentclass[10pt,a4paper]{amsart}
\usepackage[margin=19mm]{geometry}
\usepackage{amsmath,amssymb,mathtools}
\usepackage[scr=rsfs,cal=euler]{mathalfa}
\usepackage{xfrac}
\usepackage[unicode,hidelinks,bookmarks=false]{hyperref}
\newcommand{\ie}{i.e.\ }
\newcommand{\cf}{cf.\ }
\newcommand{\resp}{respectively\ }
\newcommand{\Subsection}{\subsection}
\providecommand{\nobreakdash}{\nobreak\hbox{-}}
\setlength{\emergencystretch}{2em}
\multlinegap0pt
\usepackage{bm}
\usepackage{bbm}
\usepackage{dsfont}
\usepackage{xspace}
\usepackage{cancel}
\usepackage{mathtools}
\usepackage{amsthm}
\makeatletter
\@ifundefined{theorem}{\newtheorem{theorem}{Theorem}}{}
\@ifundefined{claim}{\newtheorem{claim}[theorem]{Claim}}{}
\@ifundefined{lemma}{\newtheorem{lemma}[theorem]{Lemma}}{}
\@ifundefined{proposition}{\newtheorem{proposition}[theorem]{Proposition}}{}
\makeatother
\makeatletter
\expandafter\ifx\csname proof\endcsname\relax
\newenvironment{proof}[1][Proof]{\noindent\textbf{#1.} }{\ \rule{0.5em}{0.5em}}
\fi
\makeatother
\title[No covering with nowhere dense \textsf{P}-sets in the Cohen ]{No covering with nowhere dense \textsf{P}-sets in the Cohen model}
\author{Alan Dow, Osvaldo Guzm\'an}
\date{Source: arXiv:2507.22476v2 (CC BY 4.0)}
\begin{document}
\maketitle

\noindent\emph{Source and licence.} No covering with nowhere dense \textsf{P}-sets in the Cohen model. Version arXiv:2507.22476v2 (CC BY 4.0), distributed under the Creative Commons Attribution 4.0 International licence, as recorded in the arXiv licence metadata for this exact version and shown on the abstract page https://arxiv.org/abs/2507.22476v2. Full rights notice: licenses/arxiv-2507.22476-CC-BY-4.0.txt.\par\smallskip

\noindent\textit{Editorial scope.} Counted unit: the Proposition whose source label is \texttt{Partir sucesion} in the article (source0.tex lines 668--687, source label \texttt{Partir sucesion}), delivered together with the article's own complete original proof of it (source0.tex lines 689--875, one \texttt{proof} environment of 187 lines). No mathematics is written, repaired or re-derived: statement and proof are the article's own text line for line. Required context retained uncounted: Indescomp (lines 318--322); Lema contencion Hw1 (lines 353--358). Every definition, display and label of those lines is retained unchanged. Omitted from this package and not redistributed: 1--317, 323--352, 359--667, 688--688, 876--1684. Nothing inside a retained excerpt is omitted or altered: every retained body is delivered exactly as the article's lines are written, so no omission receipt of any kind accompanies this package. Citations: the retained excerpts carry no citation command at all, so no bibliography entry of the article is transcribed and no reference is redistributed. Reference adaptation: none was needed and none was made; every reference inside the delivered unit resolves inside it, so no unresolved reference or citation is produced. Printed slips kept literal: no printed word, symbol, hypothesis or number of the counted statement or of its proof was altered. Layout and class: the article's own class is \texttt{article} with packages including amsmath, amsfonts, amssymb, graphicx; the wrapper is the lane's pinned \texttt{amsart} editorial wrapper at 10pt on A4, which restates the theorem-like environments the retained text needs and adds the article's own macro definitions that the retained text needs (none), with extra wrapper packages (bm, bbm, dsfont, xspace, cancel, mathtools). The delivered numbering is the unit's own. Assets: no retained span contains a figure environment, a graphics-inclusion command, a PDF-page inclusion, a table or a TikZ picture; no image file is copied into this package, the package declares no asset dependency and no image file is redistributed with it. Licence: version arXiv:2507.22476v2 of this article is distributed under the Creative Commons Attribution 4.0 International licence, as recorded in the arXiv licence metadata for this exact version and shown on the abstract page https://arxiv.org/abs/2507.22476v2. A full rights notice is delivered at licenses/arxiv-2507.22476-CC-BY-4.0.txt. Characters: every retained excerpt is pure ASCII, so the delivered engine, XeLaTeX with its default Latin Modern text font, reports no missing character.
\begin{proposition}
Let $\delta$ be an indecomposable ordinal. If we split $\delta$ in finitely
many parts, then one part is isomorphic (as a linear order) to $\delta.$
\label{Indescomp}
\end{proposition}

\begin{lemma}
[\textsf{CH}]Let $M,N$ be two elements of \textsf{Sub}$\left(  \kappa\right)
.$ If $\delta_{M}\leq\delta_{N},$ then \textsf{H}$\left(  \omega_{1}\right)
\cap M\subseteq$ \textsf{H}$\left(  \omega_{1}\right)  \cap N.$
\label{Lema contencion Hw1}
\end{lemma}

\begin{proposition}
[\textsf{CH}]Assume $l,n\in\omega$ such that $l>0,$ $\left\langle
M_{0},...,M_{n}\right\rangle $ is a $\delta$-increasing sequence and $A_{i}\in
M_{i}\cap\left[  \omega_{l}\right]  ^{\leq\omega}$ for every $i\leq n.$ Let
$N$ be an $l$-model such that $\delta_{n}<\delta_{N}$ and $X\in N$ be any set.
There are $L\in N\ $and $\left\{  A_{i}\left(  u\right)  \mid i\leq n\wedge
u\in2\right\}  $ such that for every $i\leq n,$ the following properties
hold:\label{Partir sucesion}

\begin{enumerate}
\item $L\in$ \textsf{Sub}$\left(  \kappa\right)  ,$ $\delta_{n}<\delta_{L}$
and $X\in L.$

\item $\left\{  A_{i}\left(  0\right)  ,A_{i}\left(  1\right)  \right\}  \in
M_{i}$ and is a partition of $A_{i}.$

\item $A_{i}\left(  0\right)  \in L$ and $A_{i}\left(  1\right)  \cap
L=\emptyset.$
\end{enumerate}
\end{proposition}

\begin{proof}
Before starting the proof, note that (thanks to the Pairing Axiom), it is
equivalent to consider a single set $X$ or finitely many of them, since any
finite collection can be coded as a single set. The proof proceeds by
induction on $l.$ We start with $l=1.$ In here, we have that each $A_{i}$
belongs to \textsf{H}$\left(  \omega_{1}\right)  .$ Since $\delta_{i}%
<\delta_{N}$ for every $i\leq n,$ we know that $A_{0},...,A_{n}\in N$ (see
Lemma \ref{Lema contencion Hw1}). Define $A_{i}\left(  0\right)  =A_{i}$ and
$A_{i}\left(  1\right)  =\emptyset$. Since $N$ is the union of an increasing
chain of submodels, we can find $L\in N$ for which $\delta_{n},A_{0}%
,...,A_{n},X\in L.$ This finishes the proof for the base case of the induction.

\qquad\qquad\qquad\ \ \ \ 

We now assume the proposition is true for $l,$ we will prove it is true for
$l+1$ as well. Given any $\xi\in\omega_{l+1},$ denote by $g_{\xi}%
:\xi\longrightarrow\omega_{l}$ the $\trianglelefteq$-least injective function.
It follows that if $\xi\in M\in$ \textsf{Sub}$\left(  \kappa\right)
,$\ then\ $g_{\xi}\in M$. Since $N$ is an $\left(  l+1\right)  $-model, we can
find a continuous increasing chain $\left\{  N_{\alpha}\mid\alpha<\delta
_{N}\right\}  \subseteq N$ that witnesses that $N$ is an $\left(  l+1\right)
$-model and $X,\delta_{n}\in N_{0}.$

\begin{claim}
There is $\eta<\delta_{N}$ such that for every $i\leq n,$ the following
holds:\medskip

\hfill%
\begin{tabular}
[c]{l}%
$%
%TCIMACRO{\dbigcup }%
%BeginExpansion
{\displaystyle\bigcup}
%EndExpansion
\left(  N_{\eta}\cap A_{i}\right)  =%
%TCIMACRO{\dbigcup }%
%BeginExpansion
{\displaystyle\bigcup}
%EndExpansion
\left(  N_{\eta+1}\cap A_{i}\right)  .$%
\end{tabular}
\qquad\qquad\ \qquad\qquad\qquad\hfill
\end{claim}

\qquad\ \ \ \qquad\ \ \ 

We proceed by contradiction. Define $c:\delta_{N}\longrightarrow n+1$ such
that $c\left(  \eta\right)  $ is the least $i$ for which $%
%TCIMACRO{\dbigcup }%
%BeginExpansion
{\displaystyle\bigcup}
%EndExpansion
\left(  N_{\eta}\cap A_{i}\right)  <%
%TCIMACRO{\dbigcup }%
%BeginExpansion
{\displaystyle\bigcup}
%EndExpansion
\left(  N_{\eta+1}\cap A_{i}\right)  .$ By Proposition \ref{Indescomp}, we can
find $i\leq n$ for which the set $F=\left\{  \eta\in\delta_{N}\mid c\left(
\eta\right)  =i\right\}  $ is order isomorphic to $\delta_{N}.$ For any
$\eta\in F,$ let $\varepsilon_{\eta}=%
%TCIMACRO{\dbigcup }%
%BeginExpansion
{\displaystyle\bigcup}
%EndExpansion
\left(  N_{\eta}\cap A_{i}\right)  $ and define $G=\left\{  \varepsilon_{\eta
}\mid\eta\in F\right\}  .$ Note that if $\eta,\xi\in F$ and $\eta<\xi,$ then
the following holds:\medskip

\hfill%
\begin{tabular}
[c]{lll}%
$\varepsilon_{\eta}$ & $=$ & $%
%TCIMACRO{\dbigcup }%
%BeginExpansion
{\displaystyle\bigcup}
%EndExpansion
\left(  N_{\eta}\cap A_{i}\right)  $\\
& $<$ & $%
%TCIMACRO{\dbigcup }%
%BeginExpansion
{\displaystyle\bigcup}
%EndExpansion
\left(  N_{\eta+1}\cap A_{i}\right)  $\\
& $\leq$ & $%
%TCIMACRO{\dbigcup }%
%BeginExpansion
{\displaystyle\bigcup}
%EndExpansion
\left(  N_{\xi}\cap A_{i}\right)  $\\
& $=$ & $\varepsilon_{\xi}$%
\end{tabular}
\qquad\qquad\ \qquad\qquad\qquad\hfill

\qquad\ \ \ \ \qquad\ \ \ \ \ 

In this way, $F$ and $G$ are isomorphic, so the order type of $G$ is
$\delta_{N}.$ Clearly $G\subseteq\overline{A_{i}}$, but this is a
contradiction, since \textsf{OT(}$\overline{A_{i}})<\delta_{i}<\delta_{N}$.
This finishes the proof of the claim.

\qquad\qquad\ \ 

Fix $\eta<\delta_{N}$ as in the claim. For each $i\leq n,$ define:\medskip

\hfill%
\begin{tabular}
[c]{lll}%
$\beta_{i}$ & $=$ & $%
%TCIMACRO{\dbigcup }%
%BeginExpansion
{\displaystyle\bigcup}
%EndExpansion
\left(  N_{\eta}\cap A_{i}\right)  +1$\\
& $=$ & $%
%TCIMACRO{\dbigcup }%
%BeginExpansion
{\displaystyle\bigcup}
%EndExpansion
\left(  N_{\eta+1}\cap A_{i}\right)  +1.$%
\end{tabular}
\qquad\qquad\qquad\ \qquad\qquad\hfill

\qquad\ \ \ \ \qquad\ \ \ \ \ 

We now argue that $\beta_{i}\in N_{\eta+1}\cap M_{i}:$ On one hand, $\beta
_{i}\in\overline{A_{i}}\subseteq M_{i},$ and on the other hand, we have that
$\beta_{i}\in\overline{N_{\eta}\cap\omega_{l+1}}\subseteq N_{\eta+1}.$ Define
$B_{i}=g_{\beta_{i}}\left[  A_{i}\cap\beta_{i}\right]  $, which belongs to
$M_{i}$ (since both $\beta_{i}$ and $A_{i}$ are in $M_{i}$). In this way, we
have the following:

\begin{enumerate}
\item $\left\langle M_{0},...,M_{n}\right\rangle $ is a $\delta$-increasing
sequence and $B_{i}\in M_{i}\cap\left[  \omega_{l}\right]  ^{\leq\omega}$.

\item $\delta_{n}<\delta_{N_{\eta+1}}.$

\item $N_{\eta_{+1}}$ is an $l$-model.

\item $X,$ $\beta_{0},...,\beta_{n},N_{\eta}\in N_{\eta_{+1}}.$
\end{enumerate}

\qquad\ \ \ 

We can now invoke the induction hypothesis and find $L\in N_{\eta_{+1}}$ and a
sequence $\left\{  B_{i}\left(  u\right)  \mid i\leq n\wedge u\in2\right\}  $
with the following properties:

\begin{enumerate}
\item $L\in$ \textsf{Sub}$\left(  \kappa\right)  ,$ $\delta_{n}<\delta_{L}.$

\item $\left\{  B_{i}\left(  0\right)  ,B_{i}\left(  1\right)  \right\}  \in
M_{i}$ and is a partition of $B_{i}.$

\item $B_{i}\left(  0\right)  \in L$ and $B_{i}\left(  1\right)  \cap
L=\emptyset.$

\item $X,$ $\beta_{0},...,\beta_{n},N_{\eta}\in L.$
\end{enumerate}

\qquad\ \ \qquad\ \ \ 

Furthermore, since $N_{\eta}\subseteq L\subseteq N_{\eta_{+1}},$ it follows
that $\beta_{i}=%
%TCIMACRO{\dbigcup }%
%BeginExpansion
{\displaystyle\bigcup}
%EndExpansion
\left(  L\cap A_{i}\right)  +1.$ We now define $A_{i}\left(  0\right)
=g_{\beta_{i}}^{-1}\left(  B_{i}\left(  0\right)  \right)  $ and $A_{i}\left(
1\right)  =A_{i}\setminus A_{i}\left(  0\right)  .$ We will now prove this are
the items we are looking for. Since $\beta_{i}$ and $B_{i}\left(  0\right)  $
are in $M_{i}\cap L,$ it follows that $A_{i}\left(  0\right)  \in M_{i}\cap
L.$ Moreover, since $A_{i}\in M_{i},\ $we conclude that$\ A_{i}\left(
1\right)  $ is also in $M_{i}.$ It remains to prove that $A_{i}\left(
1\right)  \cap L=\emptyset.$ Assume there is $\alpha\in A_{i}\left(  1\right)
\cap L.$ It follows that $\alpha<\beta_{i},$ so we apply it the function
$g_{\beta_{i}}.$ Since $\alpha\notin A_{i}\left(  0\right)  =g_{\beta_{i}%
}^{-1}\left(  B_{i}\left(  0\right)  \right)  ,$ it follows that $g_{\beta
_{1}}\left(  \alpha\right)  \notin B_{i}\left(  0\right)  ,$ which implies
that $g_{\beta_{1}}\left(  \alpha\right)  \in B_{i}\left(  1\right)  .$
Moreover, $\beta_{i},\alpha\in L,$ so $g_{\beta_{i}}\left(  \alpha\right)  \in
L\cap B_{i}\left(  1\right)  ,$ which is a contradiction, since these two sets
are disjoint.
\end{proof}

\end{document}
